\documentclass[11pt,a4paper]{article}
\usepackage[utf8]{inputenc}
\usepackage[margin=2.4cm]{geometry}
\usepackage{amsmath,amssymb}
\usepackage{booktabs}
\usepackage{xcolor}
\usepackage{graphicx}
\usepackage{listings}
\usepackage{enumitem}
\usepackage{microtype}
\usepackage[colorlinks=true,linkcolor=blue!50!black,urlcolor=blue!50!black]{hyperref}

\definecolor{codegray}{gray}{0.45}
\definecolor{codeblue}{rgb}{0.1,0.2,0.6}
\lstset{
  language=Python,
  basicstyle=\ttfamily\footnotesize,
  keywordstyle=\color{codeblue},
  commentstyle=\color{codegray}\itshape,
  stringstyle=\color{teal},
  showstringspaces=false,
  breaklines=true,
  frame=single,
  rulecolor=\color{black!25},
  numbers=left,
  numberstyle=\tiny\color{codegray},
  literate={θ}{$\theta$}1 {ξ}{$\xi$}1,
}

\newcommand{\session}[3]{%
  \title{\textbf{#1}\\[2pt]\large #2}
  \author{astro-lab \,--\, #3}
  \date{}
  \maketitle
  \vspace{-1.5em}
}
\newcommand{\todo}[1]{\textcolor{red!60!black}{\textbf{[#1]}}}
\usepackage{framed}
\newenvironment{keyidea}{%
  \def\FrameCommand{\color{blue!40!black}\vrule width 2.5pt\hspace{8pt}\color{black}}%
  \MakeFramed{\advance\hsize-\width\FrameRestore}}{\endMakeFramed}

\newcommand{\dd}{\mathrm{d}}
\newcommand{\pc}{\,\mathrm{pc}}
\newcommand{\Msun}{M_\odot}
\newcommand{\Rsun}{R_\odot}
\setlist[itemize]{itemsep=2pt,topsep=4pt}
\setlist[enumerate]{itemsep=3pt,topsep=4pt}
\setlength{\parindent}{0pt}
\setlength{\parskip}{6pt}

\begin{document}
\session{Session 3: Stellar Structure Basics}{Hydrostatic equilibrium, polytropes, the Lane--Emden equation}{reading notes}

\section*{Goals}
Derive the equations of a simple stellar model, reduce them to one dimensionless ODE, and solve it numerically for several polytropic indices, checking against analytic solutions.

\section{Equations of structure}
For a spherical star in equilibrium, with $m(r)$ the mass inside radius $r$:
\begin{align}
\frac{\dd m}{\dd r} &= 4\pi r^2\rho, \label{eq:mass}\\
\frac{\dd P}{\dd r} &= -\frac{G\,m\,\rho}{r^2}. \label{eq:hse}
\end{align}
Equation~\eqref{eq:hse} (hydrostatic equilibrium) says pressure gradient balances gravity. A real star also needs energy transport and an equation of state; a polytrope sidesteps them with a single assumption.

\section{Polytropes}
Assume $P = K\rho^{1+1/n}$, where $n$ is the polytropic index. Rewrite~\eqref{eq:hse} as $\frac{1}{\rho}\frac{\dd P}{\dd r}=-\frac{Gm}{r^2}$, multiply by $r^2$ and differentiate, then use~\eqref{eq:mass}:
\[
\frac{1}{r^2}\frac{\dd}{\dd r}\!\left(\frac{r^2}{\rho}\frac{\dd P}{\dd r}\right) = -4\pi G\rho .
\]
Write $\rho=\rho_c\theta^n$ and $r=a\xi$ with
\[
a^2=\frac{(n+1)K\rho_c^{1/n-1}}{4\pi G}.
\]
Substituting gives the \textbf{Lane--Emden equation}
\[
\frac{1}{\xi^2}\frac{\dd}{\dd\xi}\!\left(\xi^2\frac{\dd\theta}{\dd\xi}\right) = -\theta^n,
\qquad \theta(0)=1,\ \ \theta'(0)=0 .
\]
The surface is the first zero $\xi_1$ where $\theta(\xi_1)=0$.

\subsection*{Behaviour near the centre}
The equation is singular at $\xi=0$. A series gives
\[
\theta \approx 1 - \frac{\xi^2}{6} + \frac{n\,\xi^4}{120} - \cdots
\]
so a numerical solver should start at a small $\xi_0>0$ using the series values, which is exactly what the script does.

\section{Analytic checks}
\begin{center}
\begin{tabular}{@{}lll@{}}
\toprule
$n$ & Solution $\theta(\xi)$ & $\xi_1$ \\ \midrule
0 & $1-\xi^2/6$ & $\sqrt6 \approx 2.4495$ \\
1 & $\sin\xi/\xi$ & $\pi \approx 3.1416$ \\
5 & $(1+\xi^2/3)^{-1/2}$ & $\infty$ (infinite radius, finite mass) \\
\bottomrule
\end{tabular}
\end{center}
Numerical values from the repository script, which agree with the standard tabulated values:
\begin{center}
\begin{tabular}{@{}lr@{}}
\toprule
$n$ & $\xi_1$ \\ \midrule
0 & 2.44949 \\
1 & 3.14159 \\
1.5 & 3.65375 \\
3 & 6.89685 \\
4 & 14.97155 \\
\bottomrule
\end{tabular}
\end{center}
$n=1.5$ models a fully convective star or a non-relativistic degenerate gas (white dwarf); $n=3$ is Eddington's standard model for radiation-dominated, high-mass stars.

\section{Mass and central density}
Integrating~\eqref{eq:mass} with $\rho=\rho_c\theta^n$ gives
\[
M = 4\pi a^3\rho_c\left(-\xi^2\frac{\dd\theta}{\dd\xi}\right)_{\xi_1},
\qquad R = a\,\xi_1 .
\]
The mean-to-central density ratio is $\bar\rho/\rho_c = -\dfrac{3}{\xi_1}\theta'(\xi_1)$. For $n=3$ this is about $1/54$, so such stars are strongly centrally concentrated.

\section{Code walkthrough}
\lstinputlisting{../03_stellar_structure/lane_emden.py}

Notes on the implementation: the right-hand side uses $\max(\theta,0)^n$ so the solver does not take a fractional power of a small negative number; the terminal event $\theta=0$ stops the integration at the surface.

\section{Exercises}
\begin{enumerate}
\item Plot $\theta(\xi)$ for $n=0,1,1.5,3,4$ on one axes and compare with the analytic curves where they exist.
\item Compute $M$ and $R$ for an $n=3$ polytrope with the Sun's $M$ and $R$; solve for $K$ and $\rho_c$. Compare $\rho_c$ with the true solar core density ($\sim 150\ \mathrm{g\,cm^{-3}}$).
\item Why does $n=5$ have no finite surface? Show $\theta=(1+\xi^2/3)^{-1/2}$ satisfies the equation.
\item Build a Gaia HR diagram (Session 1) and mark where $n=1.5$ and $n=3$ models would sit qualitatively.
\end{enumerate}

\section*{References}
\begin{itemize}
\item Hansen, Kawaler, Trimble, \textit{Stellar Interiors}, ch.\ 7
\item Kippenhahn, Weigert, Weiss, \textit{Stellar Structure and Evolution}
\item Chandrasekhar, \textit{An Introduction to the Study of Stellar Structure}, ch.\ IV
\end{itemize}
\end{document}
